% !PRINTED 108
\weberpage{108}
%<<CONT

$\xi$ unmöglich. Ganz ähnlich kann geschlossen werden, wenn $f(a)$ positiv und $f(b)$ negativ ist und unser Satz ist also erwiesen.

Hieraus ziehen wir zwei wichtige Folgerungen: Wenn $f(\xi)$ verschwindet, so lässt sich ein Intervall
\[
 \xi-\varepsilon<x<\xi+\varepsilon
\]
so bestimmen, dass $f(x)$ ausser für $x=\xi$ in diesem Intervall nicht verschwindet. Es sind dann vier Fälle in Bezug auf die Vorzeichen von $f(x)$ in den beiden Intervallen $\xi-\varepsilon<x<\xi$ und $\xi<x<\xi+\varepsilon$ möglich, die im folgenden Schema zusammengestellt sind:
\[
\begin{array}{lcc}
&\multicolumn{2}{c}{\xi-\varepsilon<\xi<\xi+\varepsilon}\\
1.\ f(x)&+&-\\
2.\ f(x)&-&+\\
3.\ f(x)&+&+\\
4.\ f(x)&-&-.
\end{array}
\]
Ist $f^{(\nu)}(x)$ die erste unter den Derivirten von $f(x)$, die für $x=\xi$ von Null verschieden ist, so zeigt die Formel [§. 13, (2)]
\[
 f(\xi+h)=\frac{h^\nu}{\Pi(\nu)}f^{(\nu)}(\xi)+\cdots,
\]
dass diese vier Fälle durch folgende Kennzeichen unterschieden sind:
\[
\begin{array}{lll}
1.&\nu\text{ ungerade}& f^{(\nu)}(\xi)<0\\
2.&\nu\text{ ungerade}& f^{(\nu)}(\xi)>0\\
3.&\nu\text{ gerade}& f^{(\nu)}(\xi)>0\\
4.&\nu\text{ gerade}& f^{(\nu)}(\xi)<0
\end{array}
\]
Wir sagen, dass im ersten Falle $f(x)$ \gesperrt{abnehmend}, im zweiten \gesperrt{wachsend} durch Null geht. Im Falle 3. ist $f(\xi)=0$ ein \gesperrt{Minimum}, im Falle 4. ein \gesperrt{Maximum}.

Wenn insbesondere $f'(\xi)$ von Null verschieden, $\xi$ also eine einfache Wurzel ist, so ist das positive oder negative Vorzeichen von $f'(\xi)$ das Kennzeichen für das Wachsen oder Abnehmen der Function $f(x)$ beim Durchgang von $x$ durch $\xi$.

Wir beweisen mit diesen Hülfsmitteln noch zwei weitere wichtige Sätze.

\medskip
\noindent IX. \gesperrt{Die Gleichung}\par
\[
 x^n-a=0
\]
\gesperrt{hat, wenn $a$ eine positive Zahl ist, eine und nur eine positive Wurzel.}\par
\medskip
